\section{Step 100 Publish results after the last successful stage}
Fill \texttt{publish\_results}. Recover the caller’s context address from stack offset 48. Copy the privately retained width and height from offsets 16 and 24 to context offsets 48 and 56. Store the active pixel-byte result from \texttt{\$v1} at context offset 64. Leave context inputs and unused suffix bytes unchanged.

No stage call follows these stores. The preceding result check has established packing success, and the supplied return path restores the saved return address and stack pointer without clearing the active count. Preserve zero status in \texttt{\$v0}. For the worked image, the published results are two, five and forty.

Run the complete caller in both modes. It now reports that the PNG was displayed. Repeat the bad-selector input and the documented capacity boundaries: failure must retain both framebuffer and all context bytes. Equivalent register assignments, local frame layouts and word-construction sequences are valid completed solutions. Their externally visible results and permitted accesses must agree.

\paragraph{Observe} Use \texttt{s100-1}; stop before \texttt{publication\_results\_ready}. Inspect the state described above before continuing.

\Needspace{7\baselineskip}
\paragraph{Check} Which value belongs at decoder-context offset 64?
\begin{itemize}
\item[A.] The required padded destination extent.
\item[B.] The successfully packed active pixel-byte count.
\item[C.] The original file count.
\end{itemize}
